The fundamental aspects of semidefinite programming, are closely related the ones of linear programming. Whereas in the context of linear programming the basic objects of interest were \textit{real vectors,} in the context of semidefinite programming in quantum information, the objects of interest are \textit{Hermitian operators.}

\section{Primal semidefinite programs}
A semidefinite program(SDP) is a constrained optimisation problem in an \textit{operator variable}. We shall denote this operator variable as $X$. This is a \textit{Hermitian Operator} such that $X^\dagger = X$, and acts on a complex vector space.
The objective function is a real linear function in $X$, which can always be written as $\tr(AX)$ for some hermitian operator A.

In the context of an SDP, the operator $X$ is required to satisfy some linear equality and inequality constraints, of the form $\Phi_{i}(X) = B_{i}$ for $i = 1, \ldots , m$ and $\Gamma_{j}(X) \le C_{j}$ for $j = 1, \ldots , n$ where all of the $B_{i}$ and $C_{j}$ are Hermitian operators of arbitrary finite dimension, and where $\Phi_{i}(\cdot)$ and $\Gamma_{j}(\cdot)$ are linear maps that are \textit{hermiticity-preserving.}

We use the symbol $\se$ to denote an operator inequality, i.e. $A\se B$ means that the operator $B-A$ must be positive semidefinite, i.e. all of its eigenvalues are non-negative.
Formally, an SDP can be written as
\begin{align}
    \text{maximise}
\end{align}

W.L.O.G we assume that the primal SDP is a \textit{maximisation} problem.

\begin{prf}[Maximum Eigenvalue of a Hermitian operator.]{Max Value of Hermitian operator}
    A canonical problem is to find the maximum eigenvalue of a Hermitian operator $H$, where $H$ is treated as the particular instance of the general maximum eigenvalue problem.
    We can write it as the following SDP:
    \begin{align}
        \text{maximise} \quad   & \tr(H\rho)     \\
        \text{subject to} \quad & \tr(\rho) = 1, \\
                                & \rho \se 0.
    \end{align}
    The program maximises the expected value of $H$ with respect to the quantum state $\rho$. The optimal value will be achieved when $\rho$ lies in the subspace of $H$ with the maximum eigenvalue. If $H$ is non-degenerate then this corresponds to the projector onto the eigenvector with maximum eigenvalue.
\end{prf}
In general an operator $X$ that satisfies all of the constraints, is said to be \textit{primal} feasible. The set of all primal feasible LPs $\mathcal{F}$ is convex, implying that if $X_{1}$ and $X_{2}$ are two feasible operators then $X' = pX_{1} + (1-p)X_{2}$ will also be feasible, for all $0 \le p \le 1$.

We can compactly write an SDP as
$$ \alpha = \max_{ X \in \mathcal{F}} \tr(AX),$$

with the optimal primal value being $\alpha$, and the optimal primal variable being $X^{*}$. There can be infinitely many optimal variables, which collectively form a convex set.

Similar to LPs, any feasible point $X \in \mathcal{F}$ provides a lower bound, $\tr(AX) \le \alpha$. If the feasible set of an SDP happens to be the empty set, $\mathcal{F} = \emptyset$, then the SDP is said to be infeasible, and the optimal value is taken to be $\alpha = -\infty$ i.e. the smallest possible value.
We can then write feasibility SDPs as a special class of SDPs.

\begin{prf}[Hamiltonian feasibility problem.]{Hamiltonian feasibility problem}
    The problem where we have an unknown Hamiltonian $H$, and are given the average (expected) energy of a set of quantum states, $\tr(H\rho_{i}) = E_{i},$ for $i = 1, \ldots, n.$ We need to find whether such a Hamiltonian $H$ exists whose ground state energy is non-negative, consistent with this data.
    This can be written as the following feasibility SDP:

    \begin{align}
        \text{find}  \quad      & H                                              \\
        \text{subject to} \quad & \tr(H\rho_{i}) = E_{i} \quad i = 1, \ldots, n, \\
                                & H \se 0.
    \end{align}

    A simple way to convert this into a standard optimisation SDP is to allow the hamiltonian $H$ to have negative energy eigenvalues, but to maximise the smallest eigenvalue.
    We can replace the constraint $H \se 0$  with the constraint $H \se \lambda \mathbf{I}$ where $\lambda$ is a new variable. The SDP then becomes:
    \begin{align}
        \text{maximise}   & \quad \lambda                \\
        \text{subject to} & \quad \tr(H\rho_{i}) = E_{i} \\
                          & H \se \lambda \mathbf{I}
    \end{align}
    If $\lambda^{*} < 0,$ then the problem is infeasible, i.e. there is no Hamiltonian with positive ground-state energy consistent with the data.
    If $\lambda^{*} \ge 0,$ then the problem is feasible ($H^{*}$ being a solution).
\end{prf}

Similar to the $\ell_{0}$ and $\ell_{\infty}$ norm examples where a problem which did not seem like an LP was converted to an LP, there are many problems which do not appear to be SDPs because of a non-linear objective function can be written as SDPs.\footnote{Characterising which problems admit an SDP formulation is still an active area of research.}

\begin{prf}[Trace and Operator Norms as SDPs]{norm-sdp}
    For any Hermitian operator $A$, the trace and operator norms are:
    \[
        \|A\|_1 = \sum_{i}|\lambda_i|, \quad \|A\|_\infty = \max_i|\lambda_i|,
    \]
    where $\{\lambda_i\}$ are the eigenvalues of $A$. These can each be written as an SDP:
    \begin{align}
        \|A\|_1
        = \text{minimise} \quad & \tr(X)          \\
        \text{subject to} \quad & -X \pe A \pe X,
    \end{align}
    and
    \begin{align}
        \|A\|_\infty
        = \text{minimise} \quad & x                                   \\
        \text{subject to} \quad & -x\mathbf{I} \pe A \pe x\mathbf{I},
    \end{align}
    where $x \in \mathbb{R}$ is a scalar.
\end{prf}

\begin{center}
    {\color{quantumdeep}
        \rule{\linewidth}{0.6pt}

        \scshape\Large End of Day 5

        \rule{\linewidth}{0.6pt}}
\end{center}


